\documentclass{beamer}
\usepackage[utf8]{inputenc}
\usetheme{Madrid}
\usecolortheme{default}
\usepackage{amsmath,amssymb,amsfonts,amsthm}
\usepackage{txfonts}
\usepackage{tkz-euclide}
\usepackage{listings}
\usepackage{adjustbox}
\usepackage{array}
\usepackage{tabularx}
\usepackage{gvv}
\usepackage[T1]{fontenc}
\usepackage{lmodern} % Ensure lmodern is loaded for pdflatex
\usepackage{circuitikz}
\usepackage{tikz}
\usepackage{graphicx}

\setbeamertemplate{page number in head/foot}[totalframenumber]

\usepackage{tcolorbox}
\tcbuselibrary{minted,breakable,xparse,skins}



\definecolor{bg}{gray}{0.95}
\DeclareTCBListing{mintedbox}{O{}m!O{}}{%
  breakable=true,
  listing engine=minted,
  listing only,
  minted language=#2,
  minted style=default,
  minted options={%
    linenos,
    gobble=0,
    breaklines=true,
    breakafter=,,
    fontsize=\small,
    numbersep=8pt,
    #1},
  boxsep=0pt,
  left skip=0pt,
  right skip=0pt,
  left=25pt,
  right=0pt,
  top=3pt,
  bottom=3pt,
  arc=5pt,
  leftrule=0pt,
  rightrule=0pt,
  bottomrule=2pt,
  toprule=2pt,
  colback=bg,
  colframe=orange!70,
  enhanced,
  overlay={%
    \begin{tcbclipinterior}
    \fill[orange!20!white] (frame.south west) rectangle ([xshift=20pt]frame.north west);
    \end{tcbclipinterior}},
  #3,
}
\lstset{
    language=C,
    basicstyle=\ttfamily\small,
    keywordstyle=\color{blue},
    stringstyle=\color{orange},
    commentstyle=\color{green!60!black},
    numbers=left,
    numberstyle=\tiny\color{gray},
    breaklines=true,
    showstringspaces=false,
}
%------------------------------------------------------------
%This block of code defines the information to appear in the
%Title page
\title %optional
{4.8.22}

%\subtitle{A short story}

\author % (optional)
{RATHLAVATH JEEVAN -AI25BTECH11026}



\begin{document}


\frame{\titlepage}
\begin{frame}{Question}
The distance of the line 
\begin{align}
\vec{r} = (\hat{i} - \hat{j}) + \lambda(\hat{i} + 5\hat{j} + \hat{k})
\end{align}
from the plane 
\begin{align}
\vec{r} \cdot (\hat{i} - \hat{j} + 4\hat{k}) = 5
\end{align}
is \, ?
\end{frame}
 \begin{frame}{Theoretical Solution} 
Let the line be given as:
\begin{align}
\vec{r} = \myvec{1 \\ -1 \\ 0} + \lambda \myvec{1 \\ 5 \\ 1} \label{1}
\end{align}

Let the plane be:
\begin{align}
\vec{r}^\top \myvec{1 \\ -1 \\ 4} = 5 \label{2}
\end{align}

Let:
\begin{align}
\vec{a} = \myvec{1 \\ -1 \\ 0}, \quad \vec{m} = \myvec{1 \\ 5 \\ 1}, \quad \vec{n} = \myvec{1 \\ -1 \\ 4} \label{3}
\end{align}
\end{frame}
\begin{frame}{Theoretical Solution} 
The direction vector \( \vec{m} \) is along the line.  
The normal vector \( \vec{n} \) is normal to the plane.

\textbf{Step 2: Compute \( \vec{m}^\top \vec{n} \)}

\begin{align}
\vec{m}^\top \vec{n} = \myvec{1 & 5 & 1} \myvec{1 \\ -1 \\ 4} = 1 - 5 + 4 = 0 \label{7}
\end{align}


So, the line is **parallel to the plane** and does **not intersect**.

---

Now, we use **distance from point to plane** formula geometrically.

Since line is parallel to the plane, the perpendicular distance from **any point** on the line to the plane is the required distance.

\end{frame}
\begin{frame}{Theoretical Solution} 
Let us choose \( \vec{a} = \myvec{1 \\ -1 \\ 0} \)

Let \( d \) be the distance from point \( \vec{a} \) to the plane.

Let the distance be:
\begin{align}
d = \frac{| \vec{a}^\top \vec{n} - 5 |}{\sqrt{ \vec{n}^\top \vec{n} }} \label{9}
\end{align}



Compute \( \vec{n}^\top \vec{n} \):
\begin{align}
\vec{n}^\top \vec{n} = \myvec{1 & -1 & 4} \myvec{1 \\ -1 \\ 4} = 1 + 1 + 16 = 18 \label{10}
\end{align}
\end{frame}
\begin{frame}{Theoretical Solution} 
Therefore:
\begin{align}
d = \frac{|2 - 5|}{\sqrt{18}} = \frac{3}{\sqrt{18}} = \frac{3}{3\sqrt{2}} = \frac{1}{\sqrt{2}} \label{11}
\end{align}

\[
\boxed{d = \frac{1}{\sqrt{2}}}
\]
\end{frame}
\begin{frame}{Plot}
    \centering
    \includegraphics[width=\columnwidth, height=0.8\textheight, keepaspectratio]{beamer/figs/mat8.jpeg}     
\end{frame}
\end{document}